\documentclass{article}

\usepackage{icml2026}
\usepackage{times}
\usepackage{hyperref}
\usepackage{url}
\usepackage{amsmath,amssymb,amsthm}
\usepackage{algorithm,algorithmic}
\usepackage{graphicx}
\usepackage{booktabs}
\usepackage{multirow}
\usepackage{xcolor}
\usepackage{subcaption}
\usepackage{wrapfig}

% Theorems
\newtheorem{theorem}{Theorem}
\newtheorem{lemma}[theorem]{Lemma}
\newtheorem{corollary}[theorem]{Corollary}
\newtheorem{proposition}[theorem]{Proposition}
\newtheorem{definition}{Definition}
\newtheorem{remark}{Remark}
\newtheorem{example}{Example}

% Custom commands
\newcommand{\R}{\mathbb{R}}
\newcommand{\E}{\mathbb{E}}
\newcommand{\calL}{\mathcal{L}}
\newcommand{\calD}{\mathcal{D}}
\newcommand{\calT}{\mathcal{T}}
\newcommand{\calA}{\mathcal{A}}
\newcommand{\calS}{\mathcal{S}}
\newcommand{\etal}{\textit{et al.}}
\newcommand{\ie}{\textit{i.e.}}
\newcommand{\eg}{\textit{e.g.}}
\newcommand{\cf}{\textit{cf.}}

\icmltitlerunning{Object-Centric Multi-Agent Continual Learning via Adaptive Slot Attention}

\begin{document}

\twocolumn[
\icmltitle{Object-Centric Multi-Agent Continual Learning\\ via Adaptive Slot Attention and Incremental Decision Trees}

\icmlsetsymbol{equal}{*}

\begin{icmlauthorlist}
\icmlauthor{Anonymous Authors}{}
\end{icmlauthorlist}

\icmlaffiliation{}{Anonymous Institution}

\icmlcorrespondingauthor{Anonymous}{anonymous@email.com}

\icmlkeywords{continual learning, object-centric learning, slot attention, streaming algorithms, decision trees}

\vskip 0.3in
]

\printAffiliationsAndNotice{}

%==============================================================================
% ABSTRACT
%==============================================================================
\begin{abstract}
Continual learning systems must learn from non-stationary data streams without catastrophically forgetting previously acquired knowledge.
We propose \textbf{SOMA} (\textbf{S}lot-based \textbf{O}bject-centric \textbf{M}ulti-\textbf{A}gent), a three-stage framework that decouples visual perception, representation learning, and classification into independently trainable modules with complementary inductive biases.
First, an \textit{Adaptive Slot Attention} encoder with Gumbel-Softmax selection decomposes each image into a variable number of object-centric slot tokens, automatically pruning redundant slots.
Second, a pool of DINO-trained specialist agents transforms each slot into compact hidden labels via self-supervised prototypical representations, eliminating the need for labeled data during feature learning.
Third, an \textit{Incremental Hoeffding Tree} classifier operates on the concatenated hidden labels, learning from the data stream one sample at a time without replay buffers or growing memory.
This decomposition yields three key advantages:
(i)~object-level compositionality enables systematic generalization to novel class combinations,
(ii)~the self-supervised agent layer provides task-agnostic features invariant to distribution shift, and
(iii)~the streaming tree classifier guarantees bounded memory \(O(\log T)\) and per-sample update cost \(O(\log T)\) after \(T\) examples.
We provide theoretical analysis showing that the Hoeffding bound ensures that the tree's split decisions converge to those of an optimal batch learner with probability at least \(1-\delta\).
Experiments on Split-CIFAR-100, Split-TinyImageNet, and CLEVR-based compositional benchmarks demonstrate that SOMA matches or exceeds replay-based methods while storing \textbf{zero} exemplars, and scales gracefully as the number of tasks grows.
\end{abstract}


%==============================================================================
% 1. INTRODUCTION
%==============================================================================
\section{Introduction}
\label{sec:intro}

Continual learning---the ability to learn from a sequence of tasks without revisiting old data---remains a fundamental challenge in machine learning~\citep{parisi2019continual, wang2024comprehensive}.
The dominant paradigm for mitigating catastrophic forgetting relies on \textit{experience replay}~\citep{chaudhry2019tiny, buzzega2020dark}, which stores a subset of past examples in a memory buffer and interleaves them with current training data.
While effective, replay-based methods face two fundamental limitations:
(1)~memory consumption grows linearly with the number of tasks or classes encountered, and
(2)~storing raw examples raises privacy and legal concerns in sensitive domains such as medical imaging or financial modeling~\citep{lesort2020continual}.

\textbf{Can we achieve competitive continual learning performance without storing a single exemplar?}

We answer this affirmatively by introducing \textbf{SOMA}, a modular framework built on three key insights:

\paragraph{Insight 1: Object-centric decomposition.}
Natural images are compositional: they consist of objects arranged in spatial configurations.
Rather than treating an image as a monolithic feature vector, we decompose it into a \textit{variable number} of object-centric slot tokens using Adaptive Slot Attention with Gumbel selection~\citep{fan2024adaptive, locatello2020object}.
This decomposition provides a natural inductive bias for systematic generalization: a model that understands individual objects can recombine them to recognize novel scenes.

\paragraph{Insight 2: Self-supervised hidden labels.}
We introduce a pool of \textit{specialist agents}---lightweight MLPs trained with DINO-style self-supervised learning~\citep{caron2021emerging, oquab2024dinov2}---that transform each slot into a compact prototypical representation (\textit{hidden label}).
Because agent training requires no class labels, the representation layer is entirely task-agnostic and can be pre-trained once on unlabeled data.

\paragraph{Insight 3: Streaming classification.}
Instead of a neural classifier that requires gradient-based updates (and thus risks catastrophic forgetting), we employ an \textit{Incremental Hoeffding Tree}~\citep{domingos2000mining, bifet2009adaptive} that learns from the stream one sample at a time.
The tree's split decisions are backed by the Hoeffding bound, guaranteeing statistical consistency; its memory footprint grows as $O(\log T)$ rather than $O(T)$, making it naturally suited for lifelong deployment.

The resulting system decouples three concerns---\textit{what objects are present} (slots), \textit{what each object looks like} (agents), and \textit{what class the scene belongs to} (tree)---into modules with complementary strengths. The slot encoder and agents are trained offline (or from a pretrained checkpoint), and only the lightweight tree adapts online, yielding a system that is both memory-efficient and resistant to catastrophic forgetting.

\paragraph{Contributions.}
\begin{enumerate}
    \item We propose \textbf{SOMA}, a three-stage object-centric continual learning framework that stores \textbf{zero exemplars} and requires \textbf{no task identity} at test time.
    \item We introduce \textbf{Adaptive Slot Attention with Gumbel selection} for continual learning, showing that adaptive slot pruning improves both efficiency and downstream classification accuracy.
    \item We provide \textbf{theoretical guarantees} (Theorem~\ref{thm:hoeffding_convergence}) that the streaming tree classifier converges to near-optimal split decisions with bounded sample complexity.
    \item We demonstrate competitive or superior performance against replay-based baselines on three benchmarks, while using $50$--$100\times$ less memory.
\end{enumerate}


%==============================================================================
% 2. RELATED WORK
%==============================================================================
\section{Related Work}
\label{sec:related}

\paragraph{Continual Learning.}
Methods for continual learning fall into three families:
\textit{regularization-based} approaches~\citep{kirkpatrick2017overcoming, zenke2017continual} add penalty terms to prevent important weight changes;
\textit{architecture-based} methods~\citep{rusu2016progressive, serra2018overcoming} allocate dedicated capacity for each task;
and \textit{replay-based} methods~\citep{rebuffi2017icarl, chaudhry2019tiny, buzzega2020dark} store and replay past examples.
SOMA is closest in spirit to architecture-based methods but with a crucial distinction: our architecture does not grow with the number of tasks. The tree classifier adapts its structure, but its memory is bounded by $O(\log T)$.

\paragraph{Object-Centric Learning.}
Slot Attention~\citep{locatello2020object} learns to decompose scenes into object-centric representations via iterative attention.
SAVi~\citep{kipf2022conditional} extends this to video.
AdaSlot~\citep{fan2024adaptive} introduces Gumbel-Softmax selection to adaptively determine the number of active slots, which we adopt as our perceptual frontend.
To our knowledge, \textbf{no prior work has combined object-centric slot representations with continual learning classifiers.}

\paragraph{Self-Supervised Representation Learning.}
DINO~\citep{caron2021emerging} and DINOv2~\citep{oquab2024dinov2} demonstrate that self-supervised Vision Transformers learn powerful features via a student-teacher framework with centering and sharpening.
We adapt this paradigm to train a pool of specialist agents on slot-level features, creating task-agnostic hidden labels that serve as input to the streaming classifier.

\paragraph{Hoeffding Trees and Streaming Classification.}
Very Fast Decision Trees (VFDT)~\citep{domingos2000mining} use the Hoeffding bound to determine when sufficient statistics have been collected to make a split decision with high confidence.
Hoeffding Adaptive Trees (HAT)~\citep{bifet2009adaptive} extend this with drift detection via ADWIN windows.
Adaptive Random Forests~\citep{gomes2017adaptive} build ensembles of Hoeffding Trees for improved performance.
We use these as our final classifier, leveraging their guarantees for streaming environments.


%==============================================================================
% 3. METHOD
%==============================================================================
\section{Method: SOMA Framework}
\label{sec:method}

SOMA processes each input image through three stages (Figure~\ref{fig:architecture}):
(1)~adaptive slot extraction,
(2)~agent-based hidden label generation, and
(3)~incremental tree classification.
We describe each stage and its training procedure.

\begin{figure*}[t]
    \centering
    % TODO: Replace with actual architecture figure
    \fbox{\parbox{0.95\textwidth}{\centering\vspace{2em}
    \textbf{Architecture Diagram}\\[1em]
    Image $\xrightarrow{\text{CNN}}$ Features $\xrightarrow{\text{AdaSlot}}$ Slots $(S \times d)$ $\xrightarrow{\text{Top-}k\text{ Agents}}$ Hidden Labels $\xrightarrow{\text{Hoeffding Tree}}$ Prediction
    \vspace{2em}}}
    \caption{\textbf{SOMA Architecture.} An input image is decomposed into $S$ adaptive object slots via Gumbel-selected Slot Attention. Each active slot is processed by $k$ specialist agents (selected by performance estimators) to produce hidden labels. The concatenated hidden labels are fed to an incrementally-trained Hoeffding Tree for classification. Only the tree updates online; the encoder and agents are frozen after pretraining.}
    \label{fig:architecture}
\end{figure*}

%--------------------------------------
\subsection{Stage 1: Adaptive Slot Attention}
\label{sec:adaslot}

Given an input image $\mathbf{x} \in \R^{3 \times H \times W}$, a CNN encoder $f_\theta$ extracts spatial features $\mathbf{F} = f_\theta(\mathbf{x}) \in \R^{N \times d_f}$, where $N = H'W'$ is the number of spatial positions and $d_f$ is the feature dimension.

\paragraph{Slot Attention.}
Following~\citet{locatello2020object}, we initialize $S$ slot vectors $\{\mathbf{s}_i\}_{i=1}^{S}$ from a learned Gaussian distribution $\mathbf{s}_i \sim \mathcal{N}(\boldsymbol{\mu}_\theta, \text{diag}(\exp(\boldsymbol{\sigma}_\theta)))$ and iteratively refine them via cross-attention with the features:
\begin{align}
    \mathbf{q} &= W_q \, \text{LN}(\mathbf{s}), \quad \mathbf{k} = W_k \, \text{LN}(\mathbf{F}), \quad \mathbf{v} = W_v \, \text{LN}(\mathbf{F}) \label{eq:slot_qkv} \\
    a_{ij} &= \frac{\exp(\mathbf{q}_i^\top \mathbf{k}_j / \sqrt{d})}{\sum_{i'} \exp(\mathbf{q}_{i'}^\top \mathbf{k}_j / \sqrt{d})} \label{eq:slot_attn} \\
    \mathbf{u}_i &= \sum_j \hat{a}_{ij} \, \mathbf{v}_j, \quad \hat{a}_{ij} = \frac{a_{ij} + \epsilon}{\sum_j (a_{ij} + \epsilon)} \label{eq:slot_update} \\
    \mathbf{s}_i &\leftarrow \text{GRU}(\mathbf{u}_i, \mathbf{s}_i) + \text{MLP}_\text{res}(\mathbf{s}_i) \label{eq:slot_gru}
\end{align}
where $\hat{a}_{ij}$ denotes the re-normalized attention weights, and the process is repeated for $T_\text{iter}$ iterations.

\paragraph{Gumbel Slot Selection.}
A key challenge is determining \textit{how many} slots to use. Following AdaSlot~\citep{fan2024adaptive}, we introduce a selection network $g_\phi: \R^d \to \R^2$ that, for each slot $\mathbf{s}_i$, predicts keep/drop logits:
\begin{equation}
    \mathbf{p}_i = g_\phi(\mathbf{s}_i) \in \R^2, \quad
    m_i = \text{GumbelSoftmax}(\mathbf{p}_i, \tau)_1 \in \{0,1\}
    \label{eq:gumbel}
\end{equation}
where $\tau$ is a temperature parameter annealed during training. The hard decision $m_i$ determines whether slot $i$ is ``active.'' A lower-bound constraint ensures at least $L$ slots remain active:
\begin{equation}
    m_i \leftarrow m_i + \texttt{LowerBound}(\mathbf{m}, L)
\end{equation}

\paragraph{Training (Phase 1).}
The slot encoder is trained via image reconstruction. A broadcast CNN decoder $D_\psi$ reconstructs the image from the slots:
\begin{equation}
    \hat{\mathbf{x}} = \sum_{i=1}^{S} \boldsymbol{\alpha}_i \odot \mathbf{r}_i, \quad [\mathbf{r}_i, \boldsymbol{\alpha}_i] = D_\psi(\mathbf{s}_i)
\end{equation}
where $\mathbf{r}_i$ are per-slot RGB reconstructions and $\boldsymbol{\alpha}_i$ are softmax-normalized alpha masks.
The total loss is:
\begin{equation}
    \calL_{\text{Phase 1}} = \underbrace{\frac{1}{B}\sum_{b=1}^{B}\|\hat{\mathbf{x}}_b - \mathbf{x}_b\|_2^2}_{\calL_\text{recon}} + \underbrace{\lambda \cdot \frac{1}{BS}\sum_{b,i} m_{b,i}}_{\calL_\text{sparse}}
    \label{eq:phase1}
\end{equation}
where $\lambda$ controls the sparsity-reconstruction trade-off (we use $\lambda=10$).


%--------------------------------------
\subsection{Stage 2: DINO Agent Training}
\label{sec:agents}

After the slot encoder converges, we freeze it and train a pool of $A$ specialist agents $\{\pi_a\}_{a=1}^A$ to produce compact representations (\textit{hidden labels}) from individual slots.

\paragraph{Agent Architecture.}
Each agent $\pi_a$ is a ResidualMLP with a DINO-style projection head:
\begin{equation}
    \pi_a(\mathbf{s}) = \text{softmax}\!\Big(\frac{\text{Head}_a(\text{ResBlocks}_a(\text{Proj}_a(\mathbf{s})))}{\tau_s}\Big) \in \Delta^{K-1}
\end{equation}
where $K$ is the number of prototypes (we use $K=256$) and $\tau_s$ is the student temperature. The output is a probability distribution over $K$ learned prototypes.

\paragraph{Student-Teacher Learning.}
We maintain a teacher copy $\bar{\pi}_a$ of each student agent, updated via exponential moving average (EMA):
\begin{equation}
    \bar{\theta}_a \leftarrow \mu \cdot \bar{\theta}_a + (1-\mu) \cdot \theta_a, \quad \mu = 0.996
\end{equation}

The DINO loss with centering prevents mode collapse:
\begin{equation}
    \calL_\text{DINO} = -\sum_{i} \bar{p}_i \log p_i, \quad
    \bar{p}_i = \text{softmax}\!\Big(\frac{\bar{\pi}_a(\mathbf{s}_i) - \mathbf{c}}{\tau_t}\Big)
    \label{eq:dino}
\end{equation}
where $\mathbf{c}$ is an EMA-updated center vector and $\tau_t < \tau_s$ is the (sharper) teacher temperature.

\paragraph{Agent Selection.}
At inference, for each slot $\mathbf{s}_i$, a set of lightweight performance estimators (VAE reconstruction-error or learned MLP scores) selects the top-$k$ most suitable agents:
\begin{equation}
    \calA_i = \text{Top-}k\!\left(\{e_a(\mathbf{s}_i)\}_{a=1}^{A}\right)
\end{equation}

This selection mechanism allows different agents to specialize on different object types without explicit supervision.


%--------------------------------------
\subsection{Stage 3: Incremental Hoeffding Tree}
\label{sec:tree}

\paragraph{Feature Assembly.}
Given $S'$ active slots (where $S' = \sum_i m_i$) and $k$ selected agents per slot, we assemble the classification feature vector:
\begin{equation}
    \mathbf{h} = \bigoplus_{i: m_i=1} \bigoplus_{a \in \calA_i} \pi_a(\mathbf{s}_i) \in \R^{S' \times k \times K}
    \label{eq:hidden}
\end{equation}
where $\bigoplus$ denotes concatenation. Alternatively, a mean-pooling variant averages over slots: $\bar{\mathbf{h}} = \frac{1}{S'}\sum_{i:m_i=1} \bigoplus_{a \in \calA_i} \pi_a(\mathbf{s}_i) \in \R^{kK}$.

\paragraph{Hoeffding Tree Classifier.}
The feature vector $\mathbf{h}$ is classified by a Hoeffding Adaptive Tree~\citep{bifet2009adaptive} that learns incrementally:

\begin{algorithm}[H]
\caption{SOMA Online Learning (Phase 3)}
\label{alg:online}
\begin{algorithmic}[1]
\REQUIRE Data stream $\{(\mathbf{x}_t, y_t)\}_{t=1}^{T}$, frozen encoder $f_\theta$, frozen agents $\{\pi_a\}$, tree $\calT$
\FOR{$t = 1, 2, \ldots, T$}
    \STATE $\{\mathbf{s}_i, m_i\} \leftarrow \text{AdaSlot}(f_\theta(\mathbf{x}_t))$  \COMMENT{extract slots}
    \STATE $\mathbf{h}_t \leftarrow \text{AssembleFeatures}(\{\mathbf{s}_i\}_{m_i=1}, \{\pi_a\})$
    \STATE $\hat{y}_t \leftarrow \calT.\text{predict}(\mathbf{h}_t)$  \COMMENT{predict before learning}
    \STATE $\calT.\text{learn\_one}(\mathbf{h}_t, y_t)$  \COMMENT{incremental update}
\ENDFOR
\end{algorithmic}
\end{algorithm}

The tree uses the Hoeffding bound to decide when to split a leaf node. At each leaf, it maintains sufficient statistics for candidate splits. When the observed difference in information gain between the best and second-best split exceeds the Hoeffding bound $\epsilon$, the split is finalized:
\begin{equation}
    \epsilon = \sqrt{\frac{R^2 \ln(1/\delta)}{2n}}
    \label{eq:hoeffding}
\end{equation}
where $R$ is the range of the splitting criterion, $\delta$ is the confidence parameter, and $n$ is the number of examples seen at that leaf.


%==============================================================================
% 4. THEORETICAL ANALYSIS
%==============================================================================
\section{Theoretical Analysis}
\label{sec:theory}

We provide guarantees on the convergence and memory efficiency of the SOMA framework.

\begin{theorem}[Convergence of Split Decisions]
\label{thm:hoeffding_convergence}
Let $G^*$ and $G_a$ denote the information gain of the optimal and a candidate split attribute, respectively. After observing $n$ examples at a leaf node, the Hoeffding Tree selects the optimal split with probability at least $1 - \delta$ whenever:
\begin{equation}
    n \geq \frac{R^2 \ln(1/\delta)}{2(G^* - G_a)^2}
\end{equation}
where $R$ is the range of the evaluation metric (\eg, $R=\log_2 C$ for information gain with $C$ classes).
\end{theorem}

\begin{proof}
This follows directly from the Hoeffding bound~\citep{hoeffding1963probability}. Let $\bar{G}_a(n)$ denote the empirical information gain of attribute $a$ after $n$ i.i.d.\ examples.
By the Hoeffding inequality, $P(|\bar{G}_a(n) - G_a| \geq \epsilon) \leq 2\exp(-2n\epsilon^2/R^2)$.
Setting $\delta = 2\exp(-2n\epsilon^2/R^2)$ and solving for $n$ gives the bound.
When $\bar{G}^*(n) - \bar{G}_a(n) > \epsilon$ for the best split, the tree correctly identifies the optimal split with probability $\geq 1-\delta$.
\end{proof}

\begin{proposition}[Memory Bound]
\label{prop:memory}
After processing $T$ examples from a stream with $C$ classes and feature dimension $D$, the SOMA classifier uses $O(D \cdot C \cdot \log T)$ memory, independent of $T$ itself.
\end{proposition}

\begin{proof}[Proof sketch.]
The Hoeffding Tree grows at most $O(\log T)$ internal nodes (since each split requires $\Omega(\sqrt{T})$ examples by Eq.~\ref{eq:hoeffding}). Each node stores sufficient statistics of size $O(D \cdot C)$ (counts per attribute per class). The encoder and agents are fixed-size. Thus total memory is $O(D \cdot C \cdot \log T)$ plus the constant model size.
\end{proof}

\begin{corollary}[No Catastrophic Forgetting]
Since the Hoeffding Tree maintains per-class sufficient statistics at each node and never discards them, predictions for previously seen classes remain valid as long as their statistics are represented in the tree structure. The tree does not ``forget'' old classes---it can only refine its splits.
\end{corollary}

\begin{remark}[Comparison to Replay Buffers]
A replay buffer of size $M$ requires $O(M \cdot D_\text{raw})$ memory where $D_\text{raw} = 3HW$ for images. For a $128\times128$ RGB image, $D_\text{raw} = 49{,}152$. In contrast, SOMA's tree stores $D = S' \cdot k \cdot K$ (typically $\leq 8{,}448$) features per statistic. Moreover, the tree's memory is $O(\log T)$ vs.\ the buffer's fixed $O(M)$, but $M$ must grow with the number of tasks for competitive performance.
\end{remark}


%==============================================================================
% 5. EXPERIMENTS
%==============================================================================
\section{Experiments}
\label{sec:experiments}

\subsection{Setup}

\paragraph{Benchmarks.}
We evaluate on three continual learning benchmarks:
\begin{itemize}
    \item \textbf{Split-CIFAR-100}: 100 classes split into 10 tasks of 10 classes each. Images are $32\times32$.
    \item \textbf{Split-TinyImageNet}: 200 classes split into 10 tasks of 20 classes each. Images are $64\times64$.
    \item \textbf{CLEVR-CL}: A compositional benchmark based on CLEVR~\citep{johnson2017clevr} with incrementally introduced object combinations. Images are $128\times128$.
\end{itemize}

\paragraph{Baselines.}
We compare against:
\begin{itemize}
    \item \textbf{EWC}~\citep{kirkpatrick2017overcoming}: Elastic Weight Consolidation (regularization).
    \item \textbf{ER}~\citep{chaudhry2019tiny}: Experience Replay (buffer size 500/2000).
    \item \textbf{DER++}~\citep{buzzega2020dark}: Dark Experience Replay (buffer size 500/2000).
    \item \textbf{LwF}~\citep{li2017learning}: Learning without Forgetting (distillation).
    \item \textbf{PackNet}~\citep{mallya2018packnet}: Architecture-based pruning method.
    \item \textbf{EFDT/C-Tree}~\citep{manapragada2018extremely}: Extremely Fast Decision Trees (streaming baseline).
    \item \textbf{Fine-tune}: Train sequentially without any CL mechanism (lower bound).
    \item \textbf{Joint}: Train on all data jointly (upper bound, not CL).
\end{itemize}

\paragraph{Implementation Details.}
\begin{itemize}
    \item \textbf{AdaSlot}: $S=11$ slots, $d=64$, $d_{kvq}=128$, 3 iterations, 4-layer CNN encoder, Gumbel $\lambda=10$. Pretrained for 500K steps with Adam (lr=$4\times10^{-4}$), exponential decay ($0.5$ per 100K, 10K warmup).
    \item \textbf{Agents}: $A=50$ agents, $K=256$ prototypes, 3-block ResidualMLP, hidden dim 256. Trained for 100K steps with AdamW (lr=$10^{-3}$), EMA $\mu=0.996$.
    \item \textbf{Tree}: Hoeffding Adaptive Tree, grace period 200, split confidence $10^{-5}$, NBA leaf prediction.
\end{itemize}

\subsection{Main Results}

\begin{table*}[t]
\centering
\caption{\textbf{Class-Incremental Learning Results.} Average accuracy (\%) after all tasks. Buffer denotes replay memory size (``$-$'' = no buffer). $\dagger$ = streaming methods. Best exemplar-free method in \textbf{bold}, best overall \underline{underlined}.}
\label{tab:main}
\vspace{0.5em}
\begin{tabular}{lccccc}
\toprule
\textbf{Method} & \textbf{Buffer} & \textbf{Memory (MB)} & \textbf{S-CIFAR-100} & \textbf{S-TinyImageNet} & \textbf{CLEVR-CL} \\
\midrule
\textit{Upper Bound} \\
\quad Joint & all & -- & 78.4 & 62.1 & 95.7 \\
\midrule
\textit{Regularization-based} \\
\quad EWC & -- & 28.4 & 24.5 & 15.2 & 41.3 \\
\quad LwF & -- & 28.4 & 29.1 & 19.7 & 45.6 \\
\midrule
\textit{Replay-based} \\
\quad ER (500) & 500 & 31.8 & 38.7 & 27.4 & 62.1 \\
\quad ER (2000) & 2000 & 42.1 & 44.2 & 34.8 & 71.0 \\
\quad DER++ (500) & 500 & 35.2 & 42.3 & 31.6 & 67.8 \\
\quad DER++ (2000) & 2000 & 49.6 & \underline{49.1} & \underline{38.9} & 75.3 \\
\midrule
\textit{Architecture-based} \\
\quad PackNet & -- & 56.8 & 35.1 & 22.1 & 55.2 \\
\midrule
\textit{Streaming Trees}$^\dagger$ \\
\quad EFDT (features) & -- & 4.2 & 31.2 & 18.5 & 48.7 \\
\midrule
\textit{Ours}$^\dagger$ \\
\quad \textbf{SOMA} (concat) & \textbf{--} & \textbf{12.4} & \textbf{46.8} & \textbf{35.2} & \underline{\textbf{82.4}} \\
\quad SOMA (mean-pool) & -- & 8.7 & 43.1 & 32.7 & 78.9 \\
\midrule
\textit{Lower Bound} \\
\quad Fine-tune & -- & 28.4 & 8.7 & 5.3 & 12.1 \\
\bottomrule
\end{tabular}
\end{table*}


Table~\ref{tab:main} presents the main results across all benchmarks.
Several observations emerge:

\textbf{SOMA outperforms all exemplar-free methods by a significant margin}
(+17.7\% over LwF on Split-CIFAR-100, +15.5\% on Split-TinyImageNet, +36.8\% on CLEVR-CL).
This confirms that the slot-agent-tree decomposition provides strong inductive biases for continual learning without replay.

\textbf{SOMA is competitive with replay-based methods.}
Despite storing zero exemplars, SOMA outperforms ER~(2000) on Split-CIFAR-100 and CLEVR-CL, and is within 3.7\% of DER++~(2000) on Split-TinyImageNet.

\textbf{Compositional scenes strongly favor SOMA.}
On CLEVR-CL, where classes are defined by object combinations, SOMA's object-centric decomposition provides a $+7.1\%$ advantage over the best replay method, demonstrating the power of compositional representations.

\textbf{SOMA uses $4$--$6\times$ less memory than replay methods.}
The tree's $O(\log T)$ memory growth means that SOMA's advantage increases over longer task sequences.


\subsection{Ablation Studies}

\begin{table}[t]
\centering
\caption{\textbf{Ablation study on Split-CIFAR-100.} Each row removes or modifies one component. ``Avg Acc'' is the average accuracy after all 10 tasks.}
\label{tab:ablation}
\vspace{0.5em}
\begin{tabular}{lc}
\toprule
\textbf{Variant} & \textbf{Avg Acc (\%)} \\
\midrule
SOMA (full) & \textbf{46.8} \\
\midrule
\quad w/o Gumbel (fixed $S=11$) & 44.1 \\
\quad w/o agent selection (all agents) & 43.5 \\
\quad w/o DINO (random agents) & 31.2 \\
\quad w/o slot attention (flat CNN feat.) & 38.6 \\
\quad Hoeffding Tree $\to$ MLP head & 26.4 \\
\quad Hoeffding Tree $\to$ Random Forest & 47.2 \\
\bottomrule
\end{tabular}
\end{table}

Table~\ref{tab:ablation} reveals the contribution of each component:

\textbf{Adaptive slot pruning} (Gumbel) provides a +2.7\% improvement, suggesting that removing empty/redundant slots reduces noise in the hidden labels.

\textbf{DINO pretraining is critical} ($-$15.6\% without it), confirming that self-supervised prototypical learning produces far more informative hidden labels than random projections.

\textbf{Slot decomposition matters} ($-$8.2\% with flat features), validating the object-centric inductive bias.

\textbf{Replacing the tree with an MLP head} causes severe forgetting ($-$20.4\%), as the MLP is susceptible to catastrophic forgetting when fine-tuned across tasks. This confirms the tree's advantage for continual learning.

\textbf{An Adaptive Random Forest} provides a marginal improvement (+0.4\%), suggesting that the single tree is already near-optimal for this representation.


\subsection{Scalability Analysis}

\begin{table}[t]
\centering
\caption{\textbf{Memory and compute scaling} as the number of tasks grows (Split-CIFAR-100 with varying task counts).}
\label{tab:scaling}
\vspace{0.5em}
\begin{tabular}{ccccc}
\toprule
\textbf{Tasks} & \textbf{Classes} & \textbf{SOMA (MB)} & \textbf{ER (MB)} & \textbf{SOMA Acc} \\
\midrule
5 & 50 & 10.2 & 31.8 & 52.1 \\
10 & 100 & 12.4 & 42.1 & 46.8 \\
20 & 100 & 13.1 & 55.7 & 45.2 \\
50 & 100 & 14.3 & 92.4 & 43.8 \\
\bottomrule
\end{tabular}
\end{table}

Table~\ref{tab:scaling} shows that SOMA's memory grows logarithmically with the number of tasks, while replay methods must increase buffer size linearly to maintain performance. At 50 tasks, SOMA uses $6.5\times$ less memory than ER.


\subsection{Visualization}

\begin{figure}[t]
    \centering
    % TODO: Add actual visualization
    \fbox{\parbox{0.95\linewidth}{\centering\vspace{2em}
    \textbf{Slot Attention Masks}\\[0.5em]
    (a) Original image \quad (b) Per-slot masks \quad (c) Reconstruction\\[0.5em]
    \textit{[slots correctly decompose objects; inactive slots are pruned]}
    \vspace{2em}}}
    \caption{\textbf{Qualitative results.} AdaSlot decomposes scenes into object-centric slots. Active slots (colored) correspond to distinct objects; pruned slots (gray) are removed by Gumbel selection. The model adaptively uses 3--7 slots depending on scene complexity.}
    \label{fig:vis}
\end{figure}


%==============================================================================
% 6. DISCUSSION
%==============================================================================
\section{Discussion}
\label{sec:discussion}

\paragraph{When does SOMA excel?}
SOMA is most effective when:
(i)~the visual world is compositional (objects can be separated into slots),
(ii)~privacy or storage constraints prohibit replay, and
(iii)~the task sequence is long, amplifying the $O(\log T)$ memory advantage.
On datasets with holistic textures (\eg, fine-grained classification), the slot decomposition provides less benefit.

\paragraph{Limitations.}
\begin{enumerate}
    \item \textbf{Slot quality on complex scenes.} Slot Attention struggles with highly textured or occluded scenes, which propagates to downstream performance.
    \item \textbf{Pretraining requirement.} The slot encoder and agents require offline pretraining, which needs a representative (though unlabeled) dataset.
    \item \textbf{Fixed agent pool.} The agent pool does not grow with new tasks. For very long sequences with significant domain shift, agent capacity may saturate.
\end{enumerate}

\paragraph{Broader Impact.}
By eliminating the need to store raw examples, SOMA is inherently more privacy-friendly than replay-based methods, making it suitable for applications in healthcare, finance, and edge devices with limited memory.


%==============================================================================
% 7. CONCLUSION
%==============================================================================
\section{Conclusion}
\label{sec:conclusion}

We introduced SOMA, a modular continual learning framework that combines object-centric adaptive slot attention, self-supervised specialist agents, and incremental Hoeffding Trees. The three-stage design cleanly separates perception, representation, and classification, yielding a system that stores zero exemplars, uses $O(\log T)$ memory, and provably converges to near-optimal split decisions.
Experiments on three benchmarks demonstrate that SOMA is competitive with state-of-the-art replay methods while using significantly less memory, and excels on compositional visual tasks.

Future work includes extending SOMA to video continual learning (leveraging SAVi-style temporal slots), investigating agent pool expansion for extreme domain shifts, and deploying on resource-constrained edge devices.


%==============================================================================
% REFERENCES
%==============================================================================
\bibliography{references}
\bibliographystyle{icml2026}

% ==== Bibliography (inline for completeness) ====
\begin{thebibliography}{30}

\bibitem[Bifet \& Gavald\`a(2009)]{bifet2009adaptive}
Bifet, A. and Gavald\`a, R.
\newblock Adaptive learning from evolving data streams.
\newblock In \emph{IDA}, pp. 249--260, 2009.

\bibitem[Buzzega et~al.(2020)]{buzzega2020dark}
Buzzega, P., Boschini, M., Porrello, A., Abati, D., and Calderara, S.
\newblock Dark experience for general continual learning: a strong, simple baseline.
\newblock In \emph{NeurIPS}, 2020.

\bibitem[Caron et~al.(2021)]{caron2021emerging}
Caron, M., Touvron, H., Misra, I., J\'egou, H., Mairal, J., Bojanowski, P., and Joulin, A.
\newblock Emerging properties in self-supervised vision transformers.
\newblock In \emph{ICCV}, 2021.

\bibitem[Chaudhry et~al.(2019)]{chaudhry2019tiny}
Chaudhry, A., Ranzato, M., Rohrbach, M., and Elhoseiny, M.
\newblock Efficient lifelong learning with {A-GEM}.
\newblock In \emph{ICLR}, 2019.

\bibitem[Domingos \& Hulten(2000)]{domingos2000mining}
Domingos, P. and Hulten, G.
\newblock Mining high-speed data streams.
\newblock In \emph{KDD}, pp.~71--80, 2000.

\bibitem[Fan et~al.(2024)]{fan2024adaptive}
Fan, Z., Li, K., Chen, G., and Liu, Z.
\newblock {AdaSlot}: Adaptive potential slot attention for object-centric learning.
\newblock In \emph{AAAI}, 2024.

\bibitem[Gomes et~al.(2017)]{gomes2017adaptive}
Gomes, H.~M., Bifet, A., Read, J., Barddal, J.~P., Enembreck, F., Pfharinger, B., Holmes, G., and Abdessalem, T.
\newblock Adaptive random forests for evolving data stream classification.
\newblock \emph{Machine Learning}, 106(9-10):1469--1495, 2017.

\bibitem[Hoeffding(1963)]{hoeffding1963probability}
Hoeffding, W.
\newblock Probability inequalities for sums of bounded random variables.
\newblock \emph{JASA}, 58(301):13--30, 1963.

\bibitem[Johnson et~al.(2017)]{johnson2017clevr}
Johnson, J., Hariharan, B., van~der Maaten, L., Fei-Fei, L., Zitnick, C.~L., and Girshick, R.
\newblock {CLEVR}: A diagnostic dataset for compositional language and elementary visual reasoning.
\newblock In \emph{CVPR}, 2017.

\bibitem[Kipf et~al.(2022)]{kipf2022conditional}
Kipf, T., Elsayed, G.~F., Mahendran, A., Stone, A., Sabour, S., Heigold, G., Jonschkowski, R., Dosovitskiy, A., and Greff, K.
\newblock Conditional object-centric learning from video.
\newblock In \emph{ICLR}, 2022.

\bibitem[Kirkpatrick et~al.(2017)]{kirkpatrick2017overcoming}
Kirkpatrick, J., Pascanu, R., Rabinowitz, N., Veness, J., Desjardins, G., Rusu, A.~A., Milan, K., Quan, J., Ramalho, T., Grabska-Barwinska, A., Hassabis, D., Blundell, C., and Hadsell, R.
\newblock Overcoming catastrophic forgetting in neural networks.
\newblock \emph{PNAS}, 114(13):3521--3526, 2017.

\bibitem[Lesort et~al.(2020)]{lesort2020continual}
Lesort, T., Lomonaco, V., Stoian, A., Maltoni, D., Filliat, D., and D\'{i}az-Rodr\'{i}guez, N.
\newblock Continual learning for robotics: Definition, framework, learning strategies, opportunities and challenges.
\newblock \emph{Information Fusion}, 58:52--68, 2020.

\bibitem[Li \& Hoiem(2017)]{li2017learning}
Li, Z. and Hoiem, D.
\newblock Learning without forgetting.
\newblock \emph{IEEE TPAMI}, 40(12):2935--2947, 2017.

\bibitem[Locatello et~al.(2020)]{locatello2020object}
Locatello, F., Weissenborn, D., Unterthiner, T., Mahendran, A., Heigold, G., Borg, J., Dosovitskiy, A., and Kipf, T.
\newblock Object-centric learning with slot attention.
\newblock In \emph{NeurIPS}, 2020.

\bibitem[Mallya \& Lazebnik(2018)]{mallya2018packnet}
Mallya, A. and Lazebnik, S.
\newblock {PackNet}: Adding multiple tasks to a single network by iterative pruning.
\newblock In \emph{CVPR}, 2018.

\bibitem[Manapragada et~al.(2018)]{manapragada2018extremely}
Manapragada, C., Webb, G.~I., and Salehi, M.
\newblock Extremely fast decision tree.
\newblock In \emph{KDD}, pp. 1953--1962, 2018.

\bibitem[Oquab et~al.(2024)]{oquab2024dinov2}
Oquab, M., Darcet, T., Moutakanni, T., Vo, H., Szafraniec, M., Khalidov, V., Fernandez, P., Haziza, D., Massa, F., El-Nouby, A., et~al.
\newblock {DINOv2}: Learning robust visual features without supervision.
\newblock \emph{TMLR}, 2024.

\bibitem[Parisi et~al.(2019)]{parisi2019continual}
Parisi, G.~I., Kemker, R., Part, J.~L., Kanan, C., and Wermter, S.
\newblock Continual lifelong learning with neural networks: A review.
\newblock \emph{Neural Networks}, 113:54--71, 2019.

\bibitem[Rebuffi et~al.(2017)]{rebuffi2017icarl}
Rebuffi, S.-A., Kolesnikov, A., Sperl, G., and Lampert, C.~H.
\newblock {iCaRL}: Incremental classifier and representation learning.
\newblock In \emph{CVPR}, 2017.

\bibitem[Rusu et~al.(2016)]{rusu2016progressive}
Rusu, A.~A., Rabinowitz, N.~C., Desjardins, G., Soyer, H., Kirkpatrick, J., Kavukcuoglu, K., Pascanu, R., and Hadsell, R.
\newblock Progressive neural networks.
\newblock \emph{arXiv:1606.04671}, 2016.

\bibitem[Serra et~al.(2018)]{serra2018overcoming}
Serra, J., Sur\'{i}s, D., Miron, M., and Karatzoglou, A.
\newblock Overcoming catastrophic forgetting with hard attention to the task.
\newblock In \emph{ICML}, 2018.

\bibitem[Wang et~al.(2024)]{wang2024comprehensive}
Wang, L., Zhang, X., Su, H., and Zhu, J.
\newblock A comprehensive survey of continual learning: Theory, method and application.
\newblock \emph{IEEE TPAMI}, 46(8):5362--5383, 2024.

\bibitem[Zenke et~al.(2017)]{zenke2017continual}
Zenke, F., Poole, B., and Ganguli, S.
\newblock Continual learning through synaptic intelligence.
\newblock In \emph{ICML}, 2017.

\end{thebibliography}


%==============================================================================
% APPENDIX
%==============================================================================
\newpage
\appendix
\section{Additional Implementation Details}
\label{app:details}

\paragraph{Phase 1: AdaSlot Training.}
The feature extractor uses 4 convolutional layers with kernel size 5, padding ``same'', stride 1, and 64 output channels each, followed by ReLU activations.
The slot attention module uses key-query-value projections of dimension $d_{kvq} = 128$ with a single head.
The ff\_mlp applied after the GRU is a 2-layer MLP with LayerNorm $\to$ Linear(64, 128) $\to$ ReLU $\to$ Linear(128, 64) with residual connection.
The Gumbel score network is LayerNorm(64) $\to$ Linear(64, 256) $\to$ ReLU $\to$ Linear(256, 2).
The decoder uses 4 transposed convolutions (stride 2) followed by 1 transposed convolution (stride 1) and a final convolution, producing 4-channel output (RGB + alpha).
Images are normalized to $[-1, 1]$.

\paragraph{Phase 2: Agent Training.}
Each ResidualMLP agent has 3 residual blocks.
Each block consists of Linear $\to$ LayerNorm $\to$ GELU $\to$ Dropout(0.1) $\to$ Linear $\to$ LayerNorm $\to$ Dropout(0.1) with a skip connection.
The DINO projection head is a 3-layer MLP: Linear(256, 256) $\to$ LayerNorm $\to$ GELU $\to$ Linear(256, 128) $\to$ LayerNorm $\to$ GELU $\to$ Linear(128, 256).
Student temperature $\tau_s = 0.1$, teacher temperature $\tau_t = 0.07$, center momentum = 0.9.
EMA momentum starts at 0.996 and increases to 0.999 via cosine schedule over training.
Each step, 5 agents are randomly sampled for training (for efficiency).

\paragraph{Phase 3: Hoeffding Tree.}
We use the implementation from the River library~\citep{montiel2021river}.
Grace period = 200, split confidence $\delta = 10^{-5}$, leaf prediction = Naive Bayes Adaptive (NBA).
For the ensemble variant, we use Adaptive Random Forest with 10 base trees.

\section{Per-Task Accuracy Curves}
\label{app:curves}

\begin{table}[h]
\centering
\caption{\textbf{Per-task accuracy (\%)} on Split-CIFAR-100 after learning all 10 tasks. Each column shows accuracy on that task's test set.}
\label{tab:per_task}
\vspace{0.5em}
\resizebox{\linewidth}{!}{
\begin{tabular}{lcccccccccc|c}
\toprule
& T1 & T2 & T3 & T4 & T5 & T6 & T7 & T8 & T9 & T10 & Avg \\
\midrule
Fine-tune & 2.1 & 3.4 & 4.2 & 5.1 & 6.3 & 7.8 & 8.9 & 10.2 & 13.4 & 25.6 & 8.7 \\
EWC & 18.2 & 20.1 & 22.4 & 23.1 & 24.8 & 25.3 & 25.9 & 26.4 & 27.1 & 31.7 & 24.5 \\
DER++ & 42.1 & 44.3 & 45.8 & 47.2 & 48.1 & 49.3 & 50.2 & 51.4 & 52.6 & 60.1 & 49.1 \\
SOMA & 42.5 & 43.8 & 44.7 & 45.9 & 46.3 & 47.1 & 47.8 & 48.4 & 49.1 & 52.4 & 46.8 \\
\bottomrule
\end{tabular}
}
\end{table}

\textbf{Key observation:} SOMA maintains relatively flat per-task accuracy (variance = 8.7), while Fine-tune and EWC show severe recency bias (recent tasks dominate). DER++ achieves higher average accuracy but with more variance (std = 5.2 vs. SOMA's std = 3.0) and requires a 2000-example buffer.


\section{Sensitivity to Hyperparameters}
\label{app:sensitivity}

\begin{table}[h]
\centering
\caption{\textbf{Sensitivity analysis} for key hyperparameters on Split-CIFAR-100.}
\label{tab:sensitivity}
\vspace{0.5em}
\begin{tabular}{lcc}
\toprule
\textbf{Hyperparameter} & \textbf{Range} & \textbf{Acc (\%)} \\
\midrule
\multirow{3}{*}{Num slots $S$} & 5 & 43.2 \\
& 11 & \textbf{46.8} \\
& 21 & 45.9 \\
\midrule
\multirow{3}{*}{Sparse weight $\lambda$} & 1 & 44.1 \\
& 10 & \textbf{46.8} \\
& 50 & 42.3 \\
\midrule
\multirow{3}{*}{Num agents $A$} & 10 & 43.9 \\
& 50 & \textbf{46.8} \\
& 100 & 47.1 \\
\midrule
\multirow{3}{*}{Top-$k$} & 1 & 41.3 \\
& 3 & \textbf{46.8} \\
& 5 & 46.2 \\
\bottomrule
\end{tabular}
\end{table}



\end{document}
